\chapter{Własne metody}
Autorskie metody powstały w oparciu o szkielet oraz koncepcje metod, które w badaniach wykazały najlepsze rezultaty. Krótko się przejrzaliśmy wynikom w \ref{text:results}. Wykorzystując wnioski płynące z analizy ich wzorców, opracowaliśmy nowatorskie podejścia, które skutecznie łączą precyzję z wydajnością. Rozwiązania te zostały starannie dostosowane do wymagań i złożoności problematyki wykrywania \textbf{fake news}, gwarantując wysoką jakość i efektywność.

\section{Pierwsza metoda (\textbf{Metoda17})}

\subsection{Pomysł na metodę}
\textbf{Metoda17} jest modelem zespołowym, który łączy predykcje trzech różnych klasyfikatorów: wielowarstwowego perceptronu (\texttt{MLPClassifier}), regresji logistycznej \newline(\texttt{LogisticRegression}) i XGBoosta (\texttt{XGBClassifier}) za pomocą miękkiego głosowania. Inspiracją dla tej metody jest \textbf{Metoda10}, która również wykorzystuje \texttt{VotingClassifier} do łączenia różnych modeli. Kluczowym ulepszeniem w \textbf{Metoda17} jest zastosowanie optymalizacji hiperparametrów (przynajmniej dla regresji logistycznej) oraz skalowania cech dla każdego z bazowych klasyfikatorów przed ich połączeniem. Celem jest wykorzystanie komplementarnych mocnych stron różnych algorytmów i poprawa ogólnej wydajności predykcyjnej poprzez dostrojenie parametrów i ujednolicenie skali danych wejściowych.

\subsection{Wejściowe parametry}
Metoda \texttt{metoda17} przyjmuje następujące parametry wejściowe:

\begin{enumerate}
	\def\labelenumi{\arabic{enumi}.}
	\item \textbf{X\_train}: Macierz danych treningowych dla n próbek, m cech. Każdy wiersz reprezentuje próbkę, a każda kolumna cechę. Wartości w macierzy są zazwyczaj liczbowe.
	\item \textbf{y\_train}: Wektor etykiet dla danych treningowych \textbf{n} próbek. Każdy element odpowiada próbce w \textbf{X\_train} i zawiera jej prawdziwą klasę, zazwyczaj 0 i 1 dla klasyfikacji binarnej, liczby całkowite dla wieloklasowej.
	\item \textbf{X\_test}: Macierz danych testowych. Struktura podobna do \textbf{X\_train}, używana do oceny modelu po treningu.
	\item \textbf{y\_test}: Wektor etykiet dla danych testowych. Prawdziwe klasy dla próbek w \textbf{X\_test}, używane do oceny wyników modelu.
\end{enumerate}

\subsection{Wyjściowe parametry}
Metoda zwraca następujące elementy:

\begin{enumerate}
	\def\labelenumi{\arabic{enumi}.}
	\item \textbf{voting\_clf\_final}: Wytrenowany model \texttt{VotingClassifier} po optymalizacji.
	\item \textbf{X\_test\_scaled}: Przeskalowane cechy zbioru testowego.
	\item \textbf{y\_test}: Wektor etykiet testowych.
\end{enumerate}

\subsection{Algorytm}
\begin{enumerate}
	\def\labelenumi{\arabic{enumi}.}
	\item
	      Przygotowanie i trenowanie modeli bazowych:
	      \begin{itemize}
		      \item
		            \textbf{Wielowarstwowy Perceptron (MLPClassifier)}: Inicjalizuje się model MLP. Tworzy się potok (\texttt{Pipeline}) ze skalowaniem cech (\texttt{StandardScaler}) i trenuje na \textbf{X\_train} i \textbf{y\_train}.
		      \item
		            \textbf{Regresja Logistyczna (LogisticRegression)}: Tworzy się potok (\texttt{Pipeline}) ze skalowaniem cech (\texttt{StandardScaler}) i regresją logistyczną. Definiuje się przestrzeń hiperparametrów i przeprowadza optymalizację za pomocą \texttt{GridSearchCV} na \textbf{X\_train} i \textbf{y\_train}. Wybiera się najlepszy model.
		      \item
		            \textbf{XGBoost (XGBClassifier)}: Inicjalizuje się model XGBoost. Tworzy się potok (\texttt{Pipeline}) ze skalowaniem cech (\texttt{StandardScaler}) i trenuje na \newline\textbf{X\_train} i \textbf{y\_train}.
	      \end{itemize}
	\item
	      Stworzenie i trenowanie modelu głosującego:
	      \begin{itemize}
		      \item Łączy się wytrenowane modele bazowe (MLP, zoptymalizowana regresja logistyczna, XGBoost) w model zespołowy (\texttt{VotingClassifier}) z miękkim głosowaniem.
		      \item Trenuje się model głosujący na \textbf{X\_train} i \textbf{y\_train}.
	      \end{itemize}
	\item
	      Przeskalowanie danych testowych:
	      \begin{itemize}
		      \item Używa się skalerów wytrenowanych na \textbf{X\_train} do przeskalowania \textbf{X\_test} dla każdego z modeli bazowych (osobno dla regresji logistycznej, MLP i XGBoost, jeśli były w potokach).
		      \item Ostatecznie, może zostać użyte jedno skalowanie dla całego \textbf{X\_test} do oceny modelu głosującego.
	      \end{itemize}
	\item
	      Zwracanie wyników: Zwraca się wytrenowany model głosujący (\textbf{voting\_clf\_final}), przeskalowane cechy zbioru testowego (\textbf{X\_test\_scaled}) oraz oryginalne etykiety testowe (\textbf{y\_test}).
\end{enumerate}

\subsection{Kod}
\begin{verbatim}
# Inicjalizacja modeli: MLP, XGBoost, regresja logistyczna
# w potoku ze skalowaniem i optymalizacją hiperparametrów.
mlp = MLPClassifier(...)
xgb = XGBClassifier(...)
pipeline_logreg = Pipeline(
[('scaler', StandardScaler()),
('log_reg', LogisticRegression(...))
])

# Definiowanie przestrzeni hiperparametrów do
# przeszukania dla regresji logistycznej.
param_grid_logreg = {...}

# Wyszukiwanie najlepszych hiperparametrów dla regresji logistycznej.
grid_search_logreg = GridSearchCV(pipeline_logreg, param_grid_logreg, ...)
grid_search_logreg.fit(X_train, y_train)
log_reg_optimized = grid_search_logreg.best_estimator_.named_steps['log_reg']
scaler_logreg = grid_search_logreg.best_estimator_.named_steps['scaler']
X_test_scaled_logreg = scaler_logreg.transform(X_test)

# Potok dla MLP ze skalowaniem.
pipeline_mlp = Pipeline([('scaler', StandardScaler()), ('mlp', mlp)])
pipeline_mlp.fit(X_train, y_train)
X_test_scaled_mlp = pipeline_mlp.named_steps['scaler'].transform(X_test)
mlp_optimized = pipeline_mlp.named_steps['mlp']

# Potok dla XGBoost ze skalowaniem.
pipeline_xgb = Pipeline([('scaler', StandardScaler()), ('xgb', xgb)])
pipeline_xgb.fit(X_train, y_train)
X_test_scaled_xgb = pipeline_xgb.named_steps['scaler'].transform(X_test)
xgb_optimized = pipeline_xgb.named_steps['xgb']

# Stworzenie modelu zespołowego (VotingClassifier) z miękkim głosowaniem.
voting_clf = VotingClassifier(estimators=[
('mlp', mlp_optimized),
('log_reg', log_reg_optimized),
('xgb', xgb_optimized)], voting='soft')

# Trenowanie modelu zespołowego.
voting_clf.fit(X_train, y_train)

# Przeskalowanie danych testowych.
X_test_scaled = StandardScaler().fit_transform(X_test)

return voting_clf, X_test_scaled, y_test
\end{verbatim}

\subsection{Właściwości analityczne}
\begin{enumerate}
	\def\labelenumi{\arabic{enumi}.}
	\item \textbf{Efektywność}: Łączenie różnych modeli i optymalizacja hiperparametrów regresji logistycznej ma na celu poprawę dokładności predykcji. Skalowanie cech zapewnia, że wszystkie algorytmy bazowe operują na danych o podobnej skali.
	\item \textbf{Wykorzystanie różnorodności modeli}: Zastosowanie różnych typów modeli (liniowy, nieliniowy oparty na sieci neuronowej, oparty na drzewach decyzyjnych) pozwala na uwzględnienie różnych wzorców w danych.
	\item \textbf{Optymalizacja}: Optymalizacja hiperparametrów regresji logistycznej za pomocą \texttt{GridSearchCV} pozwala na znalezienie lepszych ustawień dla tego konkretnego modelu.
	\item \textbf{Skalowalność}: Metoda może być stosowana do problemów z wieloma cechami i próbkami, jednak czas treningu może być znaczący ze względu na optymalizację hiperparametrów i trening wielu modeli.
	\item \textbf{Złożoność}: Złożoność obliczeniowa zależy od złożoności każdego z bazowych modeli oraz od procesu optymalizacji hiperparametrów: liczby przeszukiwanych kombinacji parametrów i liczby foldów walidacji krzyżowej.
\end{enumerate}